\title{The Era of 1-bit LLMs: All Large Language Models are in 1.58 Bits}

\begin{abstract}
The emergence of 1-bit Large Language Models (LLMs), as showcased by recent advancements like BitNet~\cite{bitnet}, signals a transformative phase in LLM development. In this study, we present a new 1-bit LLM variant, referred to as \textbf{\text{BitNet b1.58}}, wherein each parameter (or weight) is assigned one of three values from the ternary set \{-1, 0, 1\}. This model achieves parity with its full-precision Transformer LLM counterparts (i.e., FP16 or BF16) of equivalent size and training data, both in perplexity and downstream task results, while offering substantial advantages in terms of reduced latency, decreased memory consumption, greater throughput, and lower energy demands. More importantly, the 1.58-bit LLM establishes a novel scaling law and provides a framework for training the next generation of LLMs that achieve an optimal balance between effectiveness and efficiency. In addition, this approach introduces a new computational paradigm and paves the way for hardware innovations specifically tailored for 1-bit LLMs.
\end{abstract}

\section{The Era of 1-bit LLMs}

AI research has recently been marked by the exponential growth of Large Language Models (LLMs), which now display outstanding performance across diverse natural language processing benchmarks. Despite their successes, the increasing parameter counts of these models have introduced substantial deployment hurdles, particularly related to their significant energy requirements and associated environmental and financial costs.

A common mitigation strategy involves post-training quantization to generate low-bit versions of LLMs for inference~\cite{smoothquant, gptq, quip, quip_sharp}. This method lowers the bit-width of weights and activations, leading to notable reductions in both memory footprint and computation load. The industry has gradually shifted from 16-bit to even lower precision representations, including 4-bit solutions~\cite{gptq, awq}. Yet, while post-training quantization remains prevalent in industrial settings, its performance remains suboptimal compared to other strategies.

Recent advancements in 1-bit model architectures, exemplified by BitNet~\cite{bitnet}, indicate a promising approach to minimizing the computational expense of large language models (LLMs) while retaining robust performance. Standard LLMs typically operate with 16-bit floating point formats (such as FP16 or BF16), and a substantial portion of the computation is dedicated to matrix multiplication. The primary cost, therefore, arises from floating-point addition and multiplication operations. BitNet diverges from this by employing integer additions exclusively in its matrix multiplication routines, yielding significant energy reductions for LLM inference. As power often constrains compute throughput in hardware accelerators, these energy savings directly enable faster computations.

Beyond the arithmetic computations, loading model parameters from DRAM into the on-chip accelerator's memory (like SRAM) is another source of considerable expense during inference. While some efforts have attempted to increase SRAM size for higher throughput, this tends to incur prohibitive costs compared to DRAM. In contrast to their full-precision counterparts, 1-bit LLMs drastically shrink both memory footprint and bandwidth demands, thereby curbing the overhead and latency of transferring weights from DRAM and enabling more rapid and efficient inference.

In this paper, we present a notable variant among 1-bit LLMs: \textbf{\text{BitNet b1.58}}. This model restricts every parameter to one of three values: \{-1, 0, 1\}. Introducing 0 as a possible value, in addition to the original BitNet’s binary set, yields a 1.58-bit representation in the binary system. \text{BitNet b1.58} preserves the computational benefits of the initial 1-bit BitNet architecture, notably its near elimination of multiplication operations in matrix computation, which can be extensively optimized. It matches the energy efficiency of BitNet and achieves significantly higher memory efficiency, throughput, and lower latency than FP16-based LLMs. Furthermore, \text{BitNet b1.58} brings two additional benefits: it demonstrates stronger modeling capability via explicit feature selection—enabled by zero-valued weights that facilitate filtering—and this enhancement markedly boosts the expressivity and effectiveness of 1-bit LLMs. Secondly, our empirical findings confirm that \text{BitNet b1.58} achieves performance on par with full-precision (FP16) baselines, in terms of both perplexity and downstream tasks, beginning at the 3B parameter scale with matched training settings (such as model size and the quantity of training data).

\section{BitNet b1.58}

\text{BitNet b1.58} extends the BitNet framework, a Transformer variant where the standard \emph{nn.Linear} layers are substituted with \emph{BitLinear} modules. This model is trained from the beginning using weights quantized to 1.58 bits and activations at 8-bit precision. The model incorporates several enhancements beyond those in the original BitNet, which we detail below.

\paragraph{Quantization Function.} To ensure the weights take values in the set \{-1, 0, +1\}, we utilize an \emph{absmean} quantization procedure. This method initially normalizes the weight matrix by dividing by its mean absolute value, and subsequently applies rounding to map each entry to its closest integer in \{-1, 0, +1\}:

\begin{equation}
    \widetilde{W} = \mathrm{RoundClip}(\frac{W}{\gamma + \epsilon}, -1, 1),
\end{equation}

\begin{equation}
    \mathrm{RoundClip}(x, a, b) = \max(a, \min(b, \mathrm{round}(x))),
\end{equation}

\begin{equation}
    \gamma = \frac{1}{nm}\sum_{ij} |W_{ij}|.
\end{equation}

For activations, the quantization approach mirrors that of BitNet, but with a key difference: we omit pre-scaling the activations to the $[0, Q_b]$ interval before applying non-linearity. Instead, we normalize activations to the range $[-Q_b, Q_b]$ for each token, thereby eliminating the need for zero-point quantization. This adjustment simplifies the engineering pipeline and system optimization, while our experimental results indicate that the impact on overall model performance is insignificant.

\paragraph{LLaMA-alike Components.}

The architecture outlined by LLaMA~\cite{llama, llama2} has emerged as the default backbone within the open-source large language model landscape. To align with the broader open-source community, our version of \text{BitNet b1.58} employs components found in LLaMA. Concretely, the model integrates RMSNorm~\cite{rmsnorm}, SwiGLU~\cite{swiglu}, rotary embeddings~\cite{rope}, and omits all bias parameters. This choice ensures that \text{BitNet b1.58} can be efficiently used within widely adopted open-source libraries (such as Huggingface, vLLM~\cite{vllm}, and llama.cpp\footnote{\url{https://github.com/ggerganov/llama.cpp}}) with only minimal adaptation required.

\section{Results}

We conducted a comparative analysis between \text{BitNet b1.58} and our in-house FP16 \text{LLaMA LLM} reproduction, examining multiple model sizes. For experimental consistency, all models underwent pre-training on the RedPajama dataset~\cite{redpajama} for 100 billion tokens. Zero-shot evaluation spanned a suite of language understanding benchmarks, namely ARC-Easy~\cite{arc}, ARC-Challenge~\cite{arc}, Hellaswag~\cite{hellaswag}, Winogrande~\cite{winoGrande}, PIQA~\cite{piqa}, OpenbookQA~\cite{openbookqa}, and BoolQ~\cite{boolq}. Additionally, we assessed validation perplexity on WikiText2~\cite{wikitext2} and C4~\cite{c4} datasets.

The GPU memory footprint and inference latency of both \text{LLaMA LLM} and \text{BitNet b1.58} were benchmarked. These measurements utilized the FasterTransformer\footnote{\url{https://github.com/NVIDIA/FasterTransformer}} library, which is tailored for efficient LLM inference on GPUs. For \text{BitNet b1.58}, the evaluation incorporated the 2-bit kernel from Ladder~\cite{ladder}. We focused on reporting time per output token, as this represents the predominant computational cost in inference scenarios.

As detailed in Table~\ref{tab:ppl}, the comparison of perplexity and efficiency metrics demonstrates that \text{BitNet b1.58} achieves parity with full precision \text{LLaMA LLM} at the 3B parameter scale, yet attains 2.71$\times$ speedup and reduces GPU memory usage by a factor of 3.55. More specifically, the 3.9B parameter configuration of \text{BitNet b1.58} delivers 2.4$\times$ lower latency, uses 3.32$\times$ less GPU memory, and significantly exceeds the performance of \text{LLaMA LLM} 3B.

Table~\ref{tab:zero-shot} presents fine-grained zero-shot accuracy results across downstream tasks. Following the standardized \emph{lm-evaluation-harness}\footnote{\url{https://github.com/EleutherAI/lm-evaluation-harness}} protocol, our results indicate that the disparity in performance between \text{BitNet b1.58} and \text{LLaMA LLM} diminishes as model sizes scale. Notably, from 3B parameters onwards, \text{BitNet b1.58} achieves competitive results with the full precision reference. Consistent with perplexity trends, the downstream task evaluation highlights that \text{BitNet b1.58} at 3.9B parameters surpasses \text{LLaMA LLM} 3B, with the added benefit of reduced memory and latency. These findings establish \text{BitNet b1.58} as a Pareto optimal alternative to leading LLM architectures.

\paragraph{Memory and Latency}

We expanded our experiments by increasing the model scales to 7B, 13B, and 70B, and assessed the associated computational costs. As depicted in Figure~\ref{fig:latency-memory-energy}, both latency and memory usage exhibit distinct scaling behaviors, with the observed speed-up becoming more pronounced for larger models. Specifically, \text{BitNet b1.58} at 70B achieves a 4.1-fold speed improvement relative to the \text{LLaMA LLM} baseline. This acceleration primarily results from the fact that the computational load of \emph{nn.Linear} intensifies with growing model size. Memory requirements mirror these changes, as the embeddings—remaining at full precision—occupy a progressively smaller share of the total memory footprint for larger models. All latency and memory metrics were recorded using a 2-bit kernel, indicating that further reductions in cost could be attainable through additional optimizations.

\paragraph{Energy}

Additionally, we estimated the energy consumption arising from arithmetic operations for both \text{BitNet b1.58} and \text{LLaMA LLM}. Our evaluation centers on matrix multiplication, given its dominant impact on overall LLM computation costs. As shown in Figure~\ref{fig:energy}, the energy breakdown demonstrates that the bulk of computation in \text{BitNet b1.58} relies on INT8 additions, whereas \text{LLaMA LLM} employs a mix of FP16 addition and multiplication. Based on the chip energy modeling presented in~\cite{energycost, pokebnn}, matrix multiplication energy use for \text{BitNet b1.58} is reduced by a factor of 71.4 compared to \text{LLaMA LLM} for 7nm hardware. We also evaluated the total energy requirements for processing sequences of 512 tokens. Our findings reveal that increasing model size leads to progressively better energy efficiency for \text{BitNet b1.58} in contrast to the FP16-based \text{LLaMA LLM}. This improvement stems from the higher proportion of computations attributed to \emph{nn.Linear} in larger models, alongside the diminished influence of other module costs at scale.

\paragraph{Throughput}

We assess the throughput of \text{BitNet b1.58} and the 70B parameter \text{LLaMA LLM} on a pair of 80GB A100 GPUs, utilizing pipeline parallelism~\cite{gpipe} to enable execution of the large \text{LLaMA LLM} model. The batch size was scaled up to the point where the GPU memory became saturated, with sequences fixed at length 512. As detailed in Table~\ref{tab:throughput}, the 70B \text{BitNet b1.58} model accommodates batch sizes up to 11 times greater than those of \text{LLaMA LLM}, resulting in an 8.9-fold improvement in throughput.

\textbf{\text{BitNet b1.58} demonstrates a novel scaling law regarding the relationship between model performance and inference cost}. For comparison, Figures~\ref{fig:latency-memory-energy} and \ref{fig:energy} suggest the following approximate equivalencies between models using 1.58-bit quantization and those using 16-bit precision:

\begin{itemize}
    \item The 13B BitNet b1.58 model outperforms a 3B FP16 LLM in latency, memory footprint, and energy usage.
    \item The 30B BitNet b1.58 model surpasses a 7B FP16 LLM in latency, memory demands, and energy efficiency.
    \item The 70B BitNet b1.58 model exhibits better latency, memory consumption, and energy use than a 13B FP16 LLM.
\end{itemize}

\paragraph{Training with 2T Tokens}

The number of tokens seen during training plays a pivotal role for large language models. To examine how well \text{BitNet b1.58} scales with increased token counts, we trained a \text{BitNet b1.58} model using 2T tokens, employing the data preparation approach of StableLM-3B~\cite{StableLM-3B-4E1T}, currently a leading open-source 3B parameter model. We evaluated both models on a test suite comprising Winogrande~\cite{winoGrande}, PIQA~\cite{piqa}, SciQ~\cite{sciq}, LAMBADA~\cite{lambada}, and ARC-easy~\cite{arc}. Zero-shot performance metrics are provided in Table~\ref{tab:2t}. Where tasks offered both accuracy and normalized accuracy, their mean was reported. StableLM 3B results at 2T tokens are sourced from its technical documentation. Our analysis reveals that \text{BitNet b1.58} delivers the strongest performance across all evaluated tasks, suggesting that LLMs trained with 1.58-bit precision are also highly capable of generalization.

\section{Discussion and Future Work}

\textbf{1-bit Mixture-of-Experts (MoE) LLMs}

The Mixture-of-Experts (MoE) framework continues to be an efficient strategy for scaling LLMs economically. Despite substantially decreasing computational FLOPs, high memory requirements and the burden of inter-device communication hinder the practicality and scalability of MoE architectures. 1.58-bit LLMs offer a potential remedy to these limitations. By lowering memory usage, fewer hardware resources are necessary for MoE model deployment. In addition, the communication load for transmitting activations between devices is greatly alleviated. Ideally, if a whole model fits on a single chip, the inter-chip communication bottleneck would be eliminated entirely.

\textbf{Native Support of Long Sequence in LLMs}

With LLMs increasingly expected to process longer sequences, overcoming memory constraints—primarily from KV caches—has become crucial for efficient inference. \text{BitNet b1.58} makes substantial progress in this aspect by reducing activation precision from 16 bits to 8 bits, thus enabling the processing of sequences with double the context length using the same memory resources. There exists further opportunity to compress this down to 4 bits or less in 1.58-bit LLMs without information loss, an area we propose to explore in subsequent studies.

\textbf{LLMs on Edge and Mobile}

Deploying 1.58-bit LLMs could substantially advance the functionality of language models on edge and mobile hardware, where memory and compute resources are limited. The lower memory and power requirements characteristic of 1.58-bit LLMs make them suitable for these platforms, unlocking a suite of applications that were previously unfeasible. This would both expand the operational capacity of mobile and edge devices and create avenues for innovative LLM-powered applications. Furthermore, since 1.58-bit LLMs are well-adapted for CPUs—the dominant processors in edge and mobile systems—models such as \text{BitNet b1.58} can be run with enhanced efficiency, resulting in greater system performance.

\textbf{New Hardware for 1-bit LLMs}

Developments like Groq\footnote{\url{https://groq.com/}} underscore the promising potential for purpose-built hardware (e.g., LPUs) dedicated to LLM workloads. Building upon these advancements, there is both an opportunity and a need to design novel hardware and systems tuned for 1-bit LLMs, leveraging the unique computational characteristics introduced by architectures such as BitNet~\cite{bitnet}.